\clearpage
\VerdierPhysicalPageAnchor{0243}%
\begingroup\small\setlength{\parskip}{0pt}\setlength{\emergencystretch}{1.5em}
\setlength{\abovedisplayskip}{4pt plus 2pt minus 2pt}%
\setlength{\belowdisplayskip}{4pt plus 2pt minus 2pt}%

\VerdierSemanticAnchor{ch4d-4-5-4-proof-a-bcd}
\noindent $D^b(\mathcal A)$. The morphism $m$ induces a morphism between
the second canonical spectral sequences of the functors $F_1$ and $F_2$
relative to the complex $X$. The morphism induced between the initial terms
\[
m_2^{p,q}(X)\colon F_1^p\bigl(H^q(X)\bigr)
\longrightarrow F_2^p\bigl(H^q(X)\bigr)
\]
is an isomorphism by hypothesis. Moreover, the two spectral sequences are
stationary (4.5.3). It follows (Chapter~II, 4.4.5) that the morphism $m$
induces an isomorphism between the abutments of the two spectral sequences
and hence that $m(X)$ is an isomorphism. Assertions b), c), and d) are
proved similarly: it suffices to observe that, under the hypotheses made,
the second canonical spectral sequences are stationary (4.5.3).

\VerdierSemanticAnchor{ch4d-4-5-4-proof-e}
We prove e). By c), it suffices to show that every object of $\mathcal A$
is an object of $\operatorname{Is}(m)$. Let $X$ be an object of
$\mathcal A$ and $p$ an integer, and let us show that the morphism
\[
m(X[p])\colon F_1^p(X)\longrightarrow F_2^p(X)
\]
is an isomorphism. Since the functors $F_1$ and $F_2$ are stationary on the
left, there is an integer $r$ such that $F_i^q(Y)=0$ for $i=1,2$, every
object $Y$ of $\mathcal A$, and every integer $q\geq p+r-1$. Property
$\alpha$) implies that $X$ is quasi-isomorphic to a complex

\VerdierSemanticAnchor{ch4d-4-5-4-proof-e-complex-w}
\[
W^\bullet=
\cdots\longrightarrow0\longrightarrow Y\longrightarrow L^{-r+1}
\longrightarrow L^{-r+2}\longrightarrow\cdots\longrightarrow L^{-1}
\longrightarrow L^0\longrightarrow0\longrightarrow\cdots,
\]
where the components $L^i$, $-r+1\leq i\leq0$, are elements of $M$.
Denote by $L^\bullet$ the complex

\VerdierSemanticAnchor{ch4d-4-5-4-proof-e-complex-l}
\[
\cdots\longrightarrow0\longrightarrow L^{-r+1}
\longrightarrow L^{-r+2}\longrightarrow\cdots\longrightarrow L^{-1}
\longrightarrow L^0\longrightarrow0\longrightarrow\cdots.
\]
There is an exact sequence of complexes

\VerdierSemanticAnchor{ch4d-4-5-4-proof-e-exact-sequence}
\[
0\longrightarrow L^\bullet\longrightarrow W^\bullet
\longrightarrow Y[r]\longrightarrow0,
\]
and hence (1.3.2) a distinguished triangle in $D^*(\mathcal A)$:

\VerdierSemanticAnchor{ch4d-4-5-4-proof-e-triangle}
\[
\begin{tikzcd}[column sep=5.2em,row sep=2.8em]
{} & Y[r] \arrow[dl] \\
L^\bullet \arrow[r] & X \arrow[u]
\end{tikzcd}
\]
\endgroup
